\documentclass[11pt]{article}
\usepackage[margin=1in]{geometry}
\usepackage{amsmath,amssymb,amsthm,booktabs}
\usepackage[hidelinks]{hyperref}
\newtheorem{lemma}{Lemma}
\newtheorem{proposition}[lemma]{Proposition}
\newcommand{\dd}{\,\mathrm d}
\newcommand{\ii}{\mathrm i}
\newcommand{\ZL}{\frac{\zeta'}{\zeta}}
\newcommand{\ZS}{\mathcal S}
\newcommand{\PS}{\mathcal P}
\newcommand{\EP}{E_{\mathrm{pole}}}
\newcommand{\ET}{E_{\mathrm{trivial}}}
\newcommand{\EA}{E_{\mathrm{arch}}}
\newcommand{\ED}{E_{\mathrm{DS}}}
\title{Pair baseline: reconstruction of the unconditional explicit-formula lemma}
\author{Research proof draft}
\date{25 September 2026}
\begin{document}
\maketitle

\noindent\textbf{Task: proof. Assumptions: \texttt{UNCONDITIONAL}.}
This draft reconstructs \cite[Lemma 1, (2.2)]{BGSTB2024}, with locators
referring to arXiv:2306.04799v1. It does not reconstruct the subsequent
pair-correlation asymptotic. All nontrivial zeros
\(\rho=\beta+\ii\gamma=1/2+\delta_\rho+\ii\gamma\) are counted with
multiplicity, including both signs of \(\gamma\).
We never set \(\delta_\rho=0\) in the proof.
Powers of positive real numbers use their real logarithms.
Write \(\Lambda(p^k)=\log p\) for a prime \(p\) and \(k\geq1\),
and \(\Lambda(n)=0\) otherwise, including \(n=1\).
For \(X\geq1\), \(t\in\mathbb R\), put
\begin{align}
 \ZS(X,t)&=\sum_\rho
 \frac{2X^{\delta_\rho+\ii(\gamma-t)}}
 {1+((t-\gamma)+\ii\delta_\rho)^2},\label{eq:zero-sum}\\
 \PS(X,t)&=\sum_{n\geq1}\Lambda(n)n^{-1/2-\ii t}
                  \min\{n/X,X/n\}.\label{eq:prime-sum}
\end{align}
The interpretation of these series is absolute convergence, proved below.

\section{Analytic inputs and their provenance}\label{sec:analytic-inputs}

The following published inputs are used in their unconditional forms.
They are imported facts, not claims of new proofs in this document.
The ledger records their precise statements and source locators.
\begin{enumerate}
\item \texttt{ZETA-ANALYTIC-001}: the meromorphic continuation of zeta,
its simple pole at 1 with residue 1, its simple trivial zeros at
\(-2,-4,\ldots\), the location \(0<\Re\rho<1\) of its other zeros,
conjugation symmetry, and its absolutely convergent Euler product in
\(\Re s>1\). We use
\[
 \zeta(s)=\chi(s)\zeta(1-s),\qquad
 \chi(s)=2^s\pi^{s-1}\sin(\pi s/2)\Gamma(1-s).
\]
See \cite[\S\S25.2(i), 25.2(iv), 25.4, 25.10(i)]{DLMF};
the displayed functional equation is (25.4.2), and the Euler product
is (25.2.11). Differentiating the Euler product on compact subsets of
\(\Re s>1\) gives
\begin{equation}\label{eq:euler-derivative}
 -\ZL(s)=\sum_{n\geq1}\Lambda(n)n^{-s}.
\end{equation}
The differentiation is justified by the majorant
\(\sum_{n\geq2}(\log n)n^{-1-\varepsilon}<\infty\).
\item \texttt{ZETA-COUNT-001}: \cite[Corollary 14.3, p.~454]{MV2007}
gives
\(N_*(v)=\frac{v}{2\pi}\log\frac{v}{2\pi}-\frac{v}{2\pi}
+O(\log v)\) for \(v\geq4\), with multiplicity and half-weight at
\(\gamma=v\). This is the unconditional corollary; the following
Corollary 14.4 assumes RH and is not used.
For the project's inclusive count \(N(v)\), sandwich it between
\(N_*(v-1)\) and \(N_*(v+1)\). The two main terms differ by
\(O(\log v)\), so the same error estimate holds for \(N(v)\).
Likewise the number of occurrences in any interval \([v,v+1]\) is at
most \(N_*(v+2)-N_*(v-1)=O(\log(v+2))\) for large positive \(v\).
Conjugation and finiteness on compact sets give the corresponding bound
\(O(\log(|v|+2))\) on the whole real axis.
\item \texttt{ZETA-CONTOUR-001}: \cite[Lemma 12.2, p.~398]{MV2007}
supplies heights \(U_j\in[j,j+1]\), \(j\geq2\), avoiding zero
ordinates, such that
\(\ZL(\sigma+\ii U_j)=O(\log^2 U_j)\), uniformly for
\(-1\leq\sigma\leq2\). Conjugation supplies the same bound at
\(-U_j\). Lemma 12.4, pp.~399--400, gives
\(\ZL(z)=O(\log(|z|+1))\) uniformly for \(\Re z\leq-1\) and
\(\lvert z+2m\rvert\geq1/4\) for every positive integer \(m\).
The cited proofs use no RH or simplicity hypothesis.
\item \texttt{GAMMA-DIGAMMA-001}: \cite[Theorem C.1, (C.17), p.~523]{MV2007}
gives \(\psi(z)=\log z+O(1/|z|)\), uniformly in any fixed closed
sector bounded away from the negative real axis and from zero;
\(\psi=\Gamma'/\Gamma\) and \(\log z\) is principal.
In particular fix an effective constant \(C_\psi\) such that
\begin{equation}\label{eq:psi-bound}
 |\psi(z)-\log z|\leq C_\psi/|z|
 \quad (\Re z\geq3/2).
\end{equation}
The Euler--Maclaurin proof of that theorem yields an effective constant;
we do not optimize or numerically evaluate it here.
\end{enumerate}

\section{Convergence and a Mellin kernel}

\begin{lemma}[\texttt{PAIR-EF-CONVERGENCE-001}]
\label{lem:ef-convergence}
The zero series \eqref{eq:zero-sum} and prime series
\eqref{eq:prime-sum} converge absolutely and locally uniformly on
\([1,\infty)\times\mathbb R\). Their sums are continuous there.
For \(U\geq2(|t|+2)\),
\begin{equation}\label{eq:zero-tail}
 \sum_{|\gamma|>U}
 \left|\frac{2X^{\delta_\rho+\ii(\gamma-t)}}
 {1+((t-\gamma)+\ii\delta_\rho)^2}\right|
 \ll X^{1/2}\frac{\log U}{U},
\end{equation}
with an absolute effective implicit constant.
\end{lemma}
\begin{proof}
For real \(v\) and \(|\delta|<1/2\),
\begin{equation}\label{eq:denominator}
 |1+(v+\ii\delta)^2|
 \geq\Re(1+(v+\ii\delta)^2)
 =1+v^2-\delta^2\geq\tfrac34(1+v^2).
\end{equation}
Thus each zero term is at most \((8/3)X^{1/2}/(1+(t-\gamma)^2)\).
The inclusive counting bound above gives
\(\#\{\rho:|\gamma|\leq V\}\ll V\log(V+2)\) for \(V\geq2\).
If \(|\gamma|>U\geq2(|t|+2)\), then
\(|t-\gamma|\geq|\gamma|/2\). Summing over
\(2^kU<|\gamma|\leq2^{k+1}U\), \(k\geq0\), bounds the tail by
\[
 C X^{1/2}\sum_{k\geq0}
 \frac{\log(2^{k+1}U+2)}{2^kU}
 \ll X^{1/2}\frac{\log U}{U}.
\]
For \(1\leq X\leq B\), \(|t|\leq H\), this tail majorant is
uniform; the remaining zero terms are finite and continuous.
For primes, \(\min\{n/X,X/n\}\leq B/n\) on the same compact set,
so \(B(\log n)n^{-3/2}\) is a summable majorant, independent of \(t\).
Every term is continuous even at \(X=n\). The Weierstrass criterion
and dominated convergence prove the assertions. In particular no
principal-value ordering of the final zero series is required.
\end{proof}

\begin{lemma}[\texttt{PAIR-EF-MELLIN-001}]\label{lem:ef-mellin}
For \(y>0\),
\begin{equation}\label{eq:mellin}
 k(y):=\frac1{2\pi\ii}\int_{2-\ii\infty}^{2+\ii\infty}
          \frac{2y^w}{w^2-1}\dd w
 =\begin{cases}y-y^{-1},&y>1,\\0,&0<y\leq1.\end{cases}
\end{equation}
The integral converges absolutely, including at \(y=1\).
\end{lemma}
\begin{proof}
On \(\Re w=2\) the integrand is \(O_y((1+|\Im w|^2)^{-1})\).
For \(y>1\), move the line to \(\Re w=-R\), \(R\geq2\),
crossing residues \(y\) and \(-y^{-1}\) at \(1\) and \(-1\).
First send the horizontal heights to infinity with \(R\) fixed;
each horizontal integral is \(O_{y,R}(V^{-2})\). The new vertical
integral is \(O(y^{-R}/R)\), since
\(\lvert w^2-1\rvert\geq(R-1)^2+(\Im w)^2\).
Then send \(R\to\infty\). For \(0<y<1\), moving right crosses
no pole and the vertical bound is \(O(y^R/R)\).
At \(y=1\), moving right again gives \(O(1/R)\); thus the value
is zero. Equivalently the two left residues would cancel exactly.
These steps prove the formula without an endpoint half-weight.
\end{proof}

\section{The exact identity and its endpoints}

\begin{proposition}[\texttt{PAIR-EF-EXACT-001}]\label{prop:ef-exact}
Let \(s=1/2+\ii t\). For every \(X\geq1\), \(t\in\mathbb R\),
\begin{equation}\label{eq:exact}
 \ZS(X,t)=-\PS(X,t)-X^{-1}\ZL(s-1)+\EP(X,t)+\ET(X,t),
\end{equation}
where
\begin{align}
 \EP(X,t)&=\frac{2X^{1-s}}{s(2-s)}
       =\frac{2X^{1/2-\ii t}}{(1/2+\ii t)(3/2-\ii t)},
       \label{eq:pole}\\
 \ET(X,t)&=2\sum_{m\geq1}
 \frac{X^{-2m-s}}{(2m+s-1)(2m+s+1)}.\label{eq:trivial}
\end{align}
The trivial-zero series also converges absolutely and locally uniformly
on \([1,\infty)\times\mathbb R\).
\end{proposition}
\begin{proof}
\textit{Right-line integral.}
Fix first \(X>1\), \(t\in\mathbb R\), and define
\[
 K(z)=\frac{2X^{z-s}}{(z-s)^2-1},\qquad
 J=\frac1{2\pi\ii}\int_{5/2-\ii\infty}^{5/2+\ii\infty}
                    -\ZL(z)K(z)\dd z.
\]
There are no poles on this line. The Euler series gives a product
majorant \(2X^2\sum_n\Lambda(n)n^{-5/2}
/\lvert(2+\ii(v-t))^2-1\rvert\), integrable in \(v\).
Fubini therefore permits termwise integration. Set \(w=z-s\);
this is a shift of the vertical parameter by \(t\), justified by
absolute convergence. Lemma~\ref{lem:ef-mellin} yields
\begin{equation}\label{eq:right-line}
 J=\sum_{n\geq1}\Lambda(n)n^{-s}k(X/n)
  =\sum_{n\leq X}\Lambda(n)n^{-s}(X/n-n/X).
\end{equation}
The last sum is finite, and its endpoint term is zero.

\textit{Contours and bounds.}
For an integer \(M\geq1\) take the left line \(\Re z=-a\),
\(a=2M+1\), and horizontal sides at \(\pm U_j\) supplied by
\texttt{ZETA-CONTOUR-001}. Choose \(j\) large enough that
\(U_j>2(|t|+2)\). The rectangle is oriented with the right side
upward. In \(-1\leq\Re z\leq2\) the logarithmic derivative is
\(O(\log^2 U_j)\). On \(2\leq\Re z\leq5/2\) its Euler series
is bounded absolutely. On \(-a\leq\Re z\leq-1\), Lemma 12.4
applies on these horizontal sides and gives \(O_M(\log U_j)\).
Since \(|X^{z-s}|\leq X^2\) and the kernel denominator has size
\(\gg U_j^2\), both horizontal contributions satisfy
\begin{equation}\label{eq:horizontal-error}
 E_{\mathrm{height}}=O_{X,t,M}(\log^2 U_j/U_j^2)\longrightarrow0.
\end{equation}
This limit is taken with \(X,t,M\) fixed.

On the left line every trivial zero is at distance at least one.
Lemma 12.4 bounds its logarithmic derivative by
\(C\log(a+|v|+2)\), with \(C\) independent of \(M,v\).
Factoring the kernel denominator bounds it below by
\((a-1/2)^2+(v-t)^2\). Hence the whole left-line integral is at most
\begin{align}
 |E_{\mathrm{left}}|
 &\ll X^{-a-1/2}
 \int_{\mathbb R}\frac{\log(a+|v|+2)}{(a-1/2)^2+(v-t)^2}\dd v
 \notag\\
 &\ll X^{-a-1/2}\frac{\log(a+|t|+2)}{a}.
 \label{eq:left-error}
\end{align}
For the second inequality use \(v=t+(a-1/2)y\),
\(a+|v|+2\leq(a+|t|+2)(1+|y|)\), and the finiteness of
\(\int_{\mathbb R}\log(1+|y|)/(1+y^2)\dd y\).
After \(j\to\infty\), send \(M\to\infty\), with \(X,t\) fixed.
The bound \eqref{eq:left-error} tends to zero.

\textit{Residues.}
The residues of \(-\ZL(z)K(z)\) are as follows; at a nontrivial
zero its multiplicity multiplies the listed value.
\begin{center}
\begin{tabular}{ll}
\toprule
Location & Residue \\
\midrule
\(z=\rho\) & \(2X^{\rho-s}/(1-(\rho-s)^2)\) \\
\(z=1\) & \(-\EP(X,t)\) \\
\(z=-2m\) & \(-2X^{-2m-s}/((2m+s)^2-1)\) \\
\(z=s+1\) & \(-X\ZL(s+1)\) \\
\(z=s-1\) & \(X^{-1}\ZL(s-1)\) \\
\bottomrule
\end{tabular}
\end{center}
There are no coincidences: \(\Re(s+1)=3/2\) and
\(\Re(s-1)=-1/2\) avoid all zeta zeros and its pole, for every
real \(t\). Zeta is regular and nonzero at zero, so it supplies no
additional residue. Also
\(1-(\rho-s)^2=1+((t-\gamma)+\ii\delta_\rho)^2\).
The nontrivial-zero residues converge absolutely by
Lemma~\ref{lem:ef-convergence}. For the trivial residues,
\begin{equation}\label{eq:trivial-majorant}
 |(2m+s-1)(2m+s+1)|
 =\sqrt{((2m-1/2)^2+t^2)((2m+3/2)^2+t^2)}
 \geq m^2+t^2.
\end{equation}
Thus their absolute values are at most \(2/m^2\), uniformly for
\(X\geq1\) and real \(t\). This proves the claimed local uniform
convergence and justifies the residue limit as \(M\to\infty\).
The residue theorem now gives
\[
 J=\ZS-\EP-\ET-X\ZL(s+1)+X^{-1}\ZL(s-1).
\]
By \eqref{eq:euler-derivative} and \eqref{eq:right-line},
\[
 J+X\ZL(s+1)
 =-X^{-1}\sum_{n\leq X}\Lambda(n)n^{1-s}
   -X\sum_{n>X}\Lambda(n)n^{-1-s}=-\PS(X,t).
\]
This proves \eqref{eq:exact} for every \(X>1\), including prime powers.

\textit{Endpoint conventions and limit order.}
If \(X=n=p^k\), the coefficient \(X/n-n/X\) in \(J\) is zero.
Any common endpoint convention \(\theta\in[0,1]\), including
the Perron convention \(\theta=1/2\), multiplies this zero by
\(\theta\). After combining with the full Euler series, the coefficient
in \(-\PS\) is \(-1\), so the prime-power term is counted once.
The \(\min\) weight is continuous through every prime power, as are
all terms of \eqref{eq:exact}; the local majorants above justify
continuity of the sums.

Finally let \(X\downarrow1\), with \(t\) fixed or in a fixed compact
interval. The majorants of Lemma~\ref{lem:ef-convergence} and
\eqref{eq:trivial-majorant} justify dominated convergence.
The remaining terms are continuous, since \(s-1=-1/2+\ii t\) meets
no zero or pole of zeta. This proves the identity at \(X=1\).
The order used was: fixed \(X>1,t,M\); \(U_j\to\infty\);
then \(M\to\infty\); then \(X\downarrow1\).
No claim of uniformity of \eqref{eq:horizontal-error} in unbounded
\(X,t\) is needed for this exact identity.
\end{proof}

\section{Uniform estimates: the reconstructed lemma}

\begin{lemma}[\texttt{PAIR-EF-001}; BGSTB Lemma 1]\label{lem:pair-ef}
Uniformly for \(X\geq1\), \(t\in\mathbb R\),
\begin{align}
 \ZS(X,t)={}&-\PS(X,t)+X^{-1}\log(|t|+2)
       +\EA(X,t)+\ED(X,t)+\EP(X,t)+\ET(X,t),
       \label{eq:budget}\\
 |\EA(X,t)|&\leq C_{\rm arch}X^{-1},\qquad
 |\ED(X,t)|\leq C_{\rm DS}X^{-1},\notag\\
 |\EP(X,t)|&\leq\frac{4X^{1/2}}{1+t^2},\qquad
 |\ET(X,t)|\leq\frac{12X^{-5/2}}{|t|+2},\label{eq:error-bounds}
\end{align}
where the constants are effective and independent of \(X,t\).
Here \(\EP,\ET\) are the exact terms above and
\begin{align}
 \EA(X,t)&=X^{-1}\left(-\frac{\chi'}{\chi}(-1/2+\ii t)
                               -\log(|t|+2)\right),\notag\\
 \ED(X,t)&=X^{-1}\ZL(3/2-\ii t).\label{eq:arch-ds}
\end{align}
In particular this is precisely the error structure of
\cite[Lemma 1, (2.2)]{BGSTB2024} after combining the first two remainders.
\end{lemma}
\begin{proof}
Logarithmic differentiation of the functional equation on
\(w=-1/2+\ii t\), where all factors are finite and nonzero, gives
\[
 -\ZL(w)=-\frac{\chi'}{\chi}(w)+\ZL(1-w),\qquad
 -\frac{\chi'}{\chi}(w)
 =-\log(2\pi)-\frac\pi2\cot(\pi w/2)+\psi(1-w).
\]
No branch of \(\log\zeta\) is needed for these logarithmic
derivatives. On this line
\(\lvert\cot(-\pi/4+\ii\pi t/2)\rvert=1\), as follows by taking
the moduli of sine and cosine. Put \(z=3/2-\ii t\), \(r=|t|\).
The principal argument has modulus at most \(\pi/2\), and
\[
 \tfrac12\leq\frac{\sqrt{r^2+9/4}}{r+2}\leq1.
\]
Consequently \(\lvert\log z-\log(r+2)\rvert\leq\log2+\pi/2\).
Using \eqref{eq:psi-bound} proves the archimedean bound, for example with
\(C_{\rm arch}=\log(2\pi)+\pi+\log2+2C_\psi/3\).
This argument includes \(t=0\) and bounded \(t\), without an
asymptotic substitution \(\log|t|\) at zero.

Absolute convergence of the Euler series gives
\(C_{\rm DS}=\sum_{n\geq2}\Lambda(n)n^{-3/2}\leq
\sum_{n\geq2}(\log n)n^{-3/2}<\infty\).
For example the latter can be bounded effectively by its first terms
and an elementary integral; no prime number theorem is required.
Since
\(\sqrt{(t^2+1/4)(t^2+9/4)}\geq(1+t^2)/2\),
\eqref{eq:pole} gives the stated pole bound.
From \eqref{eq:trivial-majorant},
\[
 |\ET(X,t)|\leq2X^{-5/2}\sum_{m\geq1}\frac1{m^2+r^2}.
\]
If \(0\leq r\leq1\), the sum is at most
\(1+\int_1^\infty v^{-2}\dd v=2\), so twice the sum is at most
\(4\leq12/(r+2)\).
If \(r\geq1\), monotonicity gives
\(\sum_{m\geq1}(m^2+r^2)^{-1}\leq
\int_0^\infty(v^2+r^2)^{-1}\dd v=\pi/(2r)\).
Then \(\pi/r\leq3\pi/(r+2)\leq12/(r+2)\).
These estimates prove \eqref{eq:error-bounds}; all constants are
independent of the contour parameters, which have already been removed.
Combining only now the two \(X^{-1}\) errors gives
\[
 \ZS=-\PS+X^{-1}\bigl(\log(|t|+2)+O(1)\bigr)
 +O\!\left(\frac{X^{1/2}}{1+t^2}\right)
 +O\!\left(\frac{X^{-5/2}}{|t|+2}\right).
\]
\end{proof}

\section{Checks and scope of the audit}\label{sec:checks}

The zero-side denominator is at least \(3/4\) in modulus when
\(t=\gamma\); no exceptional real \(t\) is excluded.
Conjugation of the full zero sum and reindexing \(\rho\mapsto\bar\rho\)
give \(\ZS(X,-t)=\overline{\ZS(X,t)}\), and every exact term on the
right has the same symmetry. The reflection \(\rho\mapsto1-\bar\rho\)
preserves multiplicity and the strip bound used in all zero estimates;
it reindexes the sum with \(\delta_\rho\mapsto-\delta_\rho\), and
does not assert equality of individual summands. As a separate check
under assumption \texttt{RH}, the zero side reduces to
\(\sum_\gamma 2X^{\ii(\gamma-t)}/(1+(t-\gamma)^2)\), with
multiplicity. RH is not used elsewhere.

The rational kernel is exactly the subtraction in the source's argument:
for \(r=\rho-s\),
\[
 \frac1{1-r}-\frac1{-1-r}=\frac2{1-r^2}.
\]
This proof therefore establishes the combined identity directly, without
assigning absolute convergence to either separate Landau zero series.
There is no Fourier transform or mean-spacing rescaling in this lemma.

All local statements have status \texttt{proved-draft}. The imported
analytic inputs are identified above, the sum/integral exchanges are
justified within the proof, and no numerical result is used in it.
Independent mathematical review, comparison with the final journal
version, and reconstruction of the remaining pair-method lemmas remain
pending. A finite arithmetic regression script checks the elementary
kernel and endpoint algebra; it does not certify any analytic estimate.

\bibliographystyle{plain}
\bibliography{../research/literature}
\end{document}
